\section{Experiments}
\label{sec:experiments}

We evaluate performance across three recommendation domains, ablate the policy components, and examine archive dynamics. Section~\ref{sec:exp-setup} defines the setup and baseline taxonomy; Appendix~\ref{app:experiments} provides the detailed protocol. Section~\ref{sec:overall-performance} reports overall performance, Section~\ref{sec:component-ablation} assesses routing and feedback, and Section~\ref{sec:archive-analysis} analyzes revision dynamics. Each experimental unit is a complete corpus policy.

\subsection{Evaluation Setup}
\label{sec:exp-setup}

\paragraph{Datasets and task.}
We organize the recommendation evaluation into Amazon Beauty, Toys, and Sports.
Each evaluation example provides the user's preceding item sequence; the model's prediction consists of the next item's five characteristic words.
Table~\ref{tab:data-stats} reports the processed-domain statistics together with base-record and appended-record counts under the principal agent-run configurations.
With augmentation, every corpus retains the base records as an exact prefix.
In the formal comparisons, each eligible user contributes at most one appended record; the 26 Sports users unsupported by any evaluated method remain base-only and are not silently reassigned.

\paragraph{Baseline methods.}
We compare \method{} with all baseline methods included in Table~\ref{tab:overall}. Each baseline retains its method-specific data construction and recommendation objective, while \method{} constructs a user-conditioned policy over the available methods and adds feedback-guided revision. The same unaugmented base corpus is retained where applicable; externally reported rows retain their source protocols.

\paragraph{Data splits and reporting convention.}
Training data supply the corpus, augmentation inputs, and user features; validation feedback governs proposal revision and policy selection. To support timely feedback during iterative policy search, all compared methods are evaluated under a common one-epoch training horizon with the same backbone, data split, and fixed deterministic decoding protocol. Full split, budget, seed, training, and decoding details are in Appendix~\ref{app:experiments}; Table~\ref{tab:overall} reports the five-seed means for the matched local rows, while published comparators retain their source protocols.

\subsection{Overall Performance}
\label{sec:overall-performance}

We compare \method{} with all baseline methods included in Table~\ref{tab:overall}. The table includes SASRec~\citep{kang2018self}, BERT4Rec~\citep{sun2019bert4rec}, HSTU~\citep{zhai2024actions}, OneRec-Think~\citep{liu2025onerec_think}, and GRLM~\citep{zhang2026grlm}; externally reported rows retain their source protocols.

The validation-selected \method{} policy outperforms the strongest baseline method under the matched local protocol on the four core Recall/NDCG metrics at @10 and @20 in all three domains, giving 12/12 positive domain--metric comparisons. It also leads the local rows on HR@5, NDCG@5, and MRR@20 in each domain. This consistent ordering across Beauty, Toys, and Sports shows the value of coordinating user assignments and method exposure within one corpus. We next examine the routing and feedback components responsible for policy construction.

Protocol details and the first-round feedback ablation are provided in Appendix~\ref{app:experiments}.

\subsection{Component Ablations}
\label{sec:component-ablation}

We assess contextual routing and validation-feedback conditioning. The verified routing controls replace the route rule while preserving the domain's user population and method exposure counts. The feedback ablation compares first-round policies under matched planner, candidate-budget, data, and training conditions, varying access to validation-derived outcome information. Routing controls are fixed before evaluation. \textsc{Random} assigns users in a frozen hash order, \textsc{Length} sorts users by train-visible history length, and \textsc{Semantic} orders users by local overlap among item characteristic-word sets, dominant-word concentration, and a hash tie-breaker.

\begin{table*}[htbp]
\centering
\caption{Routing controls under matched domain budgets. Random, Length, and Semantic replace the contextual route while preserving user coverage, method exposure, and downstream training; \method{} retains the validation-selected policy.}
\label{tab:routing}
\small
\begin{tabular}{llccccccc}
\hline
Domain & Routing & HR@5 & NDCG@5 & R@10 & N@10 & R@20 & N@20 & MRR@20 \\
\hline
\multirow[c]{4}{*}{Beauty} & Random   & 0.0503 & 0.0355 & 0.0711 & 0.0423 & 0.0968 & 0.0488 & 0.0353 \\
& Length   & 0.0502 & 0.0347 & 0.0717 & 0.0417 & 0.0981 & 0.0484 & 0.0344 \\
& Semantic & 0.0459 & 0.0319 & 0.0685 & 0.0392 & 0.0939 & 0.0457 & 0.0321 \\
\rowcolor{methodblue}& RecDataAgent &\textbf{0.0563}&\textbf{0.0387}&\textbf{0.0800}&\textbf{0.0463}&\textbf{0.1087}&\textbf{0.0535}&\textbf{0.0380}\\
\hline
 \multirow[c]{4}{*}{Toys} & Random   & 0.0550 & 0.0376 & 0.0771 & 0.0447 & 0.1001 & 0.0505 & 0.0364 \\
& Length   & 0.0537 & 0.0376 & 0.0750 & 0.0445 & 0.0982 & 0.0502 & 0.0367 \\
& Semantic & 0.0527 & 0.0359 & 0.0747 & 0.0430 & 0.1008 & 0.0497 & 0.0352 \\
\rowcolor{methodblue}& RecDataAgent &\textbf{0.0592}&\textbf{0.0410}&\textbf{0.0828}&\textbf{0.0487}&\textbf{0.1108}&\textbf{0.0557}&\textbf{0.0401}\\
\hline
 \multirow[c]{4}{*}{Sports} & Random   & 0.0317 & 0.0215 & 0.0474 & 0.0265 & 0.0648 & 0.0309 & 0.0214 \\
& Length   & 0.0319 & 0.0217 & 0.0480 & 0.0268 & 0.0669 & 0.0316 & 0.0217 \\
& Semantic & 0.0331 & 0.0228 & 0.0480 & 0.0276 & 0.0669 & 0.0323 & 0.0226 \\
\rowcolor{methodblue}& RecDataAgent &\textbf{0.0385}&\textbf{0.0265}&\textbf{0.0569}&\textbf{0.0324}&\textbf{0.0779}&\textbf{0.0378}&\textbf{0.0261}\\
\hline
\end{tabular}
\end{table*}

\method{} is highest on all seven displayed metrics against Random, Length, and Semantic in each domain. With user populations and exposure counts controlled, these comparisons demonstrate the contribution of contextual routing to complete-corpus performance. Appendix~\ref{app:experiments} provides the corresponding protocol and feedback ablation.

\paragraph{Outcome-conditioned revision.}
The feedback intervention removes validation recommendation metrics, rankings, incumbent labels, and outcome-derived mechanism statements while retaining the planner configuration, method catalog, train-visible features, schema checks, and construction diagnostics. The paired first-round comparison fixes the search starting point, candidate budget, source data, and training and evaluation settings (Appendix~\ref{app:experiments}). Generated policies can change in response to the intervention. Table~\ref{tab:no-feedback} compares the final validation-selected policy and the no-feedback candidate under equal augmentation and per-candidate training budgets: both use the same base corpus, domain-specific appended-record count, one epoch of full-parameter fine-tuning, and evaluation protocol.

\begin{table*}[htbp]
\centering
\caption{End-to-end comparison of the validation-selected and no-feedback policies under matched corpus and training budgets. Both conditions retain the same base corpus, domain-specific appended-record count, one-epoch fine-tuning, and deterministic evaluation; Table~\ref{tab:feedback-raw} isolates feedback availability at the first round.}
\label{tab:no-feedback}
\small
\begin{tabular}{llccccccc}
\hline
Domain & Policy & HR@5 & NDCG@5 & R@10 & N@10 & R@20 & N@20 & MRR@20 \\
\hline
\rowcolor{methodblue}\multirow[c]{2}{*}{Beauty} & \method{}    &\textbf{0.0563}&\textbf{0.0387}&\textbf{0.0800}&\textbf{0.0463}&\textbf{0.1087}&\textbf{0.0535}&\textbf{0.0380}\\
       & No feedback & 0.0544 & 0.0380 & 0.0770 & 0.0453 & 0.1056 & 0.0525 & 0.0376 \\
\rowcolor{methodblue}\multirow[c]{2}{*}{Toys}   & \method{}    &\textbf{0.0592}&\textbf{0.0410}&\textbf{0.0828}&\textbf{0.0487}&\textbf{0.1108}&\textbf{0.0557}&\textbf{0.0401}\\
       & No feedback & 0.0537 & 0.0364 & 0.0756 & 0.0435 & 0.1027 & 0.0503 & 0.0355 \\
\rowcolor{methodblue}\multirow[c]{2}{*}{Sports} & \method{}    &\textbf{0.0385}&\textbf{0.0265}&\textbf{0.0569}&\textbf{0.0324}&\textbf{0.0779}&\textbf{0.0378}&\textbf{0.0261}\\
       & No feedback & 0.0308 & 0.0211 & 0.0456 & 0.0259 & 0.0631 & 0.0303 & 0.0211 \\
\hline
\end{tabular}
\end{table*}

The two comparisons answer different questions. Table~\ref{tab:no-feedback} shows that the final validation-selected policy outperforms the no-feedback candidate on all displayed metrics in each domain. Table~\ref{tab:feedback-raw} isolates feedback availability at the first round under matched conditions: feedback improves all four core metrics in Beauty and Sports, while Toys trades lower Recall for higher NDCG. The matched first-round ablation shows gains in Beauty and Sports and a Recall--NDCG trade-off in Toys; the final-policy comparison describes the outcome of the complete search procedure. We next examine successive validation-guided revisions.

\subsection{Archive Dynamics}
\label{sec:archive-analysis}

We analyze three successive validation-feedback-conditioned policies under a common protocol and include the final validation-selected policy as a separate reference. Table~\ref{tab:three-rounds} reports the complete seven-metric trajectory. Beauty recovers after Round~2, Toys is lower in Rounds~2--3, and Sports shows a Recall--NDCG trade-off in Round~2 followed by lower scores in Round~3. These fixed-sequence profiles complement the final policy and show that archive updates follow validation utility rather than proposal recency. The separate Sports batch from a second planner backend is documented in the appendix as an execution-portability trace.

\begin{table*}[b]
\centering
\caption{Validation-guided policy sequence and selected reference. Rounds 1--3 are successive policies in the revision sequence; Ref.\ is the separately selected final policy reported in Table~\ref{tab:overall}.}
\label{tab:three-rounds}
\footnotesize
\setlength{\tabcolsep}{2.8pt}
\begin{tabular}{llcccccccc}
\hline
Domain & Round & HR@5 & NDCG@5 & R@10 & N@10 & R@20 & N@20 & MRR@20 & Outcome profile \\
\hline
\multirow[c]{4}{*}{Beauty} & 1 & 0.0552 & 0.0382 & 0.0791 & 0.0459 & 0.1066 & 0.0528 & 0.0377 & Round reference \\
       & 2 & 0.0443 & 0.0293 & 0.0679 & 0.0369 & 0.0954 & 0.0439 & 0.0295 & Lower scores \\
       & 3 & 0.0562 &\textbf{0.0393}& 0.0783 &\textbf{0.0464}& 0.1075 &\textbf{0.0538}&\textbf{0.0387}& Recovery \\
 \rowcolor{methodblue}      & Ref. &\textbf{0.0563}& 0.0387 &\textbf{0.0800}& 0.0463 &\textbf{0.1087}& 0.0535 & 0.0380 & Validation-selected \\
\hline
\multirow[c]{4}{*}{Toys}   & 1 & 0.0530 & 0.0367 & 0.0749 & 0.0437 & 0.1012 & 0.0504 & 0.0361 & Round reference \\
       & 2 & 0.0521 & 0.0355 & 0.0725 & 0.0421 & 0.0956 & 0.0479 & 0.0344 & Lower scores \\
       & 3 & 0.0490 & 0.0333 & 0.0693 & 0.0399 & 0.0951 & 0.0464 & 0.0326 & Lower scores \\
 \rowcolor{methodblue}      & Ref. &\textbf{0.0592}&\textbf{0.0410}&\textbf{0.0828}&\textbf{0.0487}&\textbf{0.1108}&\textbf{0.0557}&\textbf{0.0401}& Validation-selected \\
\hline
\multirow[c]{4}{*}{Sports} & 1 & 0.0328 & 0.0220 & 0.0496 & 0.0275 & 0.0692 & 0.0324 & 0.0221 & Round reference \\
       & 2 & 0.0332 & 0.0226 & 0.0493 & 0.0278 & 0.0682 & 0.0325 & 0.0225 & Recall--NDCG trade-off \\
       & 3 & 0.0316 & 0.0216 & 0.0472 & 0.0266 & 0.0675 & 0.0318 & 0.0218 & Lower scores \\
 \rowcolor{methodblue}      & Ref. &\textbf{0.0385}&\textbf{0.0265}&\textbf{0.0569}&\textbf{0.0324}&\textbf{0.0779}&\textbf{0.0378}&\textbf{0.0261}& Validation-selected \\
\hline
\end{tabular}
\end{table*}
\FloatBarrier

